\section{Error Analysis}
\label{sec:errors}

To uncover the qualitative mechanisms behind retrieval breakdown, we conducted an error analysis on the 32 development queries across all three domains, focusing on queries where models scored $\text{nDCG}@10 = 0.0$.

\subsection{Head-to-Head Model Comparison (Dense vs. BM25)}
Evaluating performance deltas ($\Delta = \text{nDCG}_{\text{Dense}} - \text{nDCG}_{\text{BM25}}$) revealed:
\begin{itemize}
    \item \textbf{Dense Substantially Beats BM25 ($\Delta > 0.10$):} 18 queries (56.3\% of the dev set). In domains like \texttt{law} and \texttt{workplace}, the dense bi-encoder successfully recovered gold documents that shared virtually no vocabulary with the query.
    \item \textbf{BM25 Substantially Beats Dense ($\Delta < -0.10$):} 0 queries. BM25 never decisively outperformed the dense encoder on any query in this subset.
    \item \textbf{Joint Failure ($\text{nDCG} = 0.0$ for both):} 14 queries.
\end{itemize}

\subsection{Categorization of Failure Modes}
We categorized the 14 joint failure cases into three distinct error mechanisms, summarized in Table~\ref{tab:error_categories}:

\begin{table}[h]
\centering
\small
\caption{Taxonomy of failure modes on challenging development queries.}
\label{tab:error_categories}
\begin{tabular}{lcp{9.0cm}}
\toprule
\textbf{Error Category} & \textbf{Count} & \textbf{Linguistic / Operational Mechanism} \\
\midrule
\textbf{Domain / Lexical Gap} & 10 & Highly specialized jargon or formal mathematical formulas (prominent in \texttt{quant}) where neither bag-of-words nor general-purpose text bi-encoders could map colloquial user questions to formal statutory or mathematical expressions. \\
\addlinespace
\textbf{Multi-Step Reasoning} & 3 & Queries requiring multi-hop deduction (e.g., resolving a dependent legal clause or historical prerequisite prior to identifying the relevant policy). Single-pass whole-query retrieval fails without step-by-step query decomposition. \\
\addlinespace
\textbf{Temporal Mismatch} & 1 & The model successfully matched the topical semantics of the query but retrieved passages describing an outdated or subsequent year (e.g., retrieving 2021 regulatory reforms for a question targeting 2014). \\
\bottomrule
\end{tabular}
\end{table}

\subsection{Per-Domain Gap and Qualitative Insights}
\begin{itemize}
    \item \textbf{\texttt{quant} (Hardest Domain, Dense 0.0919):} Dense and sparse models both struggle severely. Financial questions frequently embed LaTeX equations, numeric variables, and specialized abbreviations (e.g., ``Black-Scholes delta hedging under jump diffusion''). BGE-small, trained on natural prose, maps these formulas to generic financial text rather than the mathematically exact gold response.
    \item \textbf{\texttt{law} (Easiest Domain, Dense 0.3697, LightGBM 0.1340):} Legal questions are rich in structured narrative context, enabling dense bi-encoders to capture topical semantics effectively. Furthermore, explicit year annotations allow the LightGBM reranker's \texttt{year\_overlap} feature to filter out temporally conflicting statutes.
    \item \textbf{\texttt{workplace} (Dense 0.2676, BM25 0.0230):} Human resources questions (e.g., severance negotiations, non-compete enforcement) are phrased conversationally by users, but answered formally by forum specialists. BM25 fails completely due to vocabulary mismatch, whereas semantic embeddings bridge the terminology gap.
\end{itemize}
